% Narrow binary packing with a complete post-exponent range argument.
\section{Narrow binary packing: 89 operations}
\label{sec:narrow-binary-product}

The binary coefficient compiler of Section~\ref{sec:binary-product}
permits a narrower packing. A change in its lowest mask digit removes
one multiplication from the power chain. The proof must recover the
mask's size after exponent decoding and then exclude a possible carry
into the fixed-code field; neither property is assumed in advance.

\begin{theorem}\label{thm:best89}
Every recursively enumerable subset of the positive integers has the
same admissible fixed index $(V,H,T_L)$ as in Theorem~\ref{thm:best90}
for a system with 34 positive unknowns and 22 equations admitting a
straight-line certificate of 89 operations: 47 multiplications and
42 additions. The supplied inputs are $(x,V,H,T_L)$, with three fixed
components besides the positive queried input $x$. Fixed numerals and
equality tests are free.
\end{theorem}

\subsection{The changed source and preliminary bounds}

Retain the compiler, fixed index, product equation, positive bounds,
geometry, and Pell equations of Section~\ref{sec:binary-product}.
In particular,
\begin{equation}\label{eq:narrow-binary-retained}
\begin{gathered}
 C=x+g,\qquad \sigma=(e-\ell)C^2,\qquad
 \ell+\sigma+\alpha=q,\\
 b=x+\beta,\qquad \theta=H+b=B-2,\qquad
 \lambda(B-1)=q^2-1,\\
 S_2=\ell+eq=V+t\theta.
\end{gathered}
\end{equation}
All supplied unknowns are strictly positive. Replace only the packing
definitions and their index equation by
\begin{equation}\label{eq:narrow-binary-packing}
\begin{aligned}
 n&=q^4, & S&=g+q\sigma+q^2S_2,\\
 T_+&=q^2\theta\lambda+\ell(\theta q-b), &
 r&=S(n^2-n)+T_+(n^2-1).
\end{aligned}
\end{equation}
Thus the unchanged bound gives
\[
 0<g,\sigma,\ell,e<q,\qquad 0<S<q^4=n.
\]
Before any exponent is decoded, geometry gives $B\le q^2$ and
$\theta\lambda>b$. Therefore
\[
 T_+>bq^2-b\ell>0,\qquad
 T_+<q^4+\theta\ell q<2q^4=2n.
\]
Here $\theta\lambda=q^2-1-\lambda<q^2$ and
$\theta\ell q<Bq^2\le q^4$. The fixed compiler threshold then gives
\begin{equation}\label{eq:narrow-binary-preliminary}
 n\ge64,\qquad n\le n^2-1\le r<3n^3,\qquad
 2\le L,\qquad 3L\le B\le n.
\end{equation}
In particular, the narrower bound $T_+<n$ has not yet been used.

\subsection{Pell recovery with the wider preliminary range}

Use the unchanged Pell definitions and equations
\eqref{eq:binary-product-pell-definitions}--\eqref{eq:binary-product-second-pell}
at this new $n$ and $r$. Their proof extends to
\eqref{eq:narrow-binary-preliminary}. Indeed,
\[
 U,Y_P\ge n^2,\qquad
 E\ge n^4>3n^3\ge r+1,\qquad
 a>n^4>6n^3>2r+1.
\]
The first Pell index is at least $r+1$, the main index is at least
$r+2$, and the relaxed auxiliary and half-parameter arguments recover
the exact main index $J=2r+1$. The same Pell growth comparison then
recovers the first index $r+1$.

For $\xi=(U+1)^{2r}/U^r$, the lower ratio bound remains
$c/k>\xi$. The only changed estimate in the upper bound is
\[
 \frac{4r}{a}<\frac{12}{n}\le\frac3{16}<\frac12.
\]
Consequently the same elementary estimate gives
\[
 \frac ck<\xi\left(1+\frac{8r}{a}\right),\qquad
 Y_P\ge U^r,\qquad a>U^{r+1},\qquad
 0<\frac ck-\xi<\frac{16r}{U+1}.
\]
The first exponent criterion proves $U=2^{2r+1}$, independently of
mask decoding. The positive gap $c>\kappa$ and fixed index congruence
give the second Pell index $L$, since $0<L<J$ and both values
$\psi_2(L),\psi_2(t_0)$ are less than $a$. Its size hypotheses are
unchanged:
\[
 B^{3L}\le B^B\le n^n\le U^r<a,
 \qquad q^3\le n^n\le U^r<a.
\]
Thus the second exponent criterion proves $q=B^L$. Both $B$ and $q$
are powers of two. Exact binomial symmetry and the ratio estimate give
the same rounding conclusion as before,
\begin{equation}\label{eq:narrow-binary-pell-output}
 n^2\mid\binom{2r}{r}.
\end{equation}
This proof applies to arbitrary positive source solutions. It has used
neither canonical fixed codes nor nonintersection of packed fields.

\subsection{Recover the packing range after exponent decoding}

Now $q=B^L$ with $L\ge2$, so
\[
 \lambda=\frac{q^2-1}{B-1}\ge B+1>\theta.
\]
Using $\ell<q$ in the exact identity
\[
 T_+=q^4-q^2-\lambda q^2+\ell(\theta q-b)
\]
shows $T_+<q^4=n$. Together with $0<S<n$, the ordinary packed
binary lemma therefore applies to \eqref{eq:narrow-binary-pell-output}:
\begin{equation}\label{eq:narrow-binary-disjoint}
 S\mathbin{\&}(T_+-1)=0.
\end{equation}
There is no appeal to that lemma outside its proved range.

\subsection{A spill into the fixed-code mask is impossible}

Write
\[
 b\ell=k_0q+r_0,\qquad 0\le r_0<q,
 \qquad \theta\ell-1-k_0=s_0q+u_0,\qquad 0\le u_0<q.
\]
Since $b<B\le q$, we have $0\le k_0\le b-1$. Also
$\theta\ell-1-k_0\ge\theta-b=H>0$, and it is less than
$\theta q$. Hence $0\le s_0\le B-3$. The normalized mask is
\begin{equation}\label{eq:narrow-binary-normalized}
 T_+-1=(q-1-r_0)+qu_0+q^2(\theta\lambda+s_0).
\end{equation}
The last coefficient is below $q^2$, because
$\theta\lambda=q^2-1-\lambda$ and $s_0<\lambda$.
The corresponding fields of $S$ are $g,\sigma,S_2$, of widths
$q,q,q^2$. Their bounds justify separate extraction from
\eqref{eq:narrow-binary-disjoint}.

The base-$B$ digits of $\theta\lambda$ are all $B-2$. If $s_0=1$,
the unit mask digit becomes $B-1$, while the others remain $B-2$.
Every digit of $S_2$, and therefore of $\ell$, is then Boolean, so
\[
 \ell\le\frac{q-1}{B-1}<\frac q{B-1},\qquad \theta\ell<q.
\]
This contradicts $s_0=1$. If $s_0\ge2$, the first two mask digits
become $s_0-2,B-1$ and all later digits remain $B-2$. The unit
digit of $\ell$ is at most $B-1$, its next digit is zero, and its
higher digits are Boolean. Since $L\ge2$,
\[
 \ell\le B-1+\sum_{i=2}^{L-1}B^i
       =\frac{q-2B+1}{B-1}<\frac q{B-1}.
\]
Again $\theta\ell<q$, contradicting $s_0\ge2$. Thus $s_0=0$.

The unperturbed middle mask now proves that $S_2$ has a Boolean
digit polynomial $F(T)$ of degree below $2L$. Reduction of
$S_2=V+t(B-2)$ modulo $B-2$ gives $F(2)\equiv V\pmod{B-2}$.
The fixed threshold \eqref{eq:binary-product-threshold} puts both
integers below $B-2$. Binary uniqueness therefore recovers
\begin{equation}\label{eq:narrow-binary-canonical}
 \ell=\ell_0(B),\qquad e=e_0(B).
\end{equation}
The unchanged support bound $K+1<L$ implies $b\ell<q$ and
$\theta\ell<q$. Consequently $k_0=0$, and the other two masks are
exactly
\begin{equation}\label{eq:narrow-binary-individual-masks}
 g\mathbin{\&}(q-1-b\ell)=0,
 \qquad \sigma\mathbin{\&}(\theta\ell-1)=0.
\end{equation}

\subsection{The altered low mask preserves every circuit test}

The first mask in \eqref{eq:narrow-binary-individual-masks} is the
same coordinate mask as before, with its unused high block removed.
At an indicator position its digit is $H_0=B-1-b$, and elsewhere
it is $B-1$. Since $g<q$, the removed high block makes no difference.
It recovers the physical coordinate polynomial $C(B)=x+g$ exactly.

The second mask differs from $\theta\ell$ only at and below the
least indicator position $v_0=8$. Below $v_0$ its digits are $B-1$;
at $v_0$ the digit is $B-3$ instead of $B-2$. But the unchanged
compiler has $\min\operatorname{supp}D>M\ge v_0$. Independently of
the values of the coordinate digits,
\[
 \sigma=D(B)(x+g)^2\equiv0\pmod{B^{v_0+1}}.
\]
Every altered digit of $\sigma$ is therefore zero. The new test is
equivalent to the old product mask. All coefficient-isolation,
carry-reset, helper-normalization and unit arguments of
Section~\ref{sec:binary-product} apply unchanged. This proves the
soundness direction for the same fixed index and the same raw input.

\subsection{Positive converse and exact operation count}

Given a represented input, retain the canonical coding witnesses of
Section~\ref{sec:binary-product} at a sufficiently large power-of-two
radix. They satisfy all unchanged bounds and give the three masks
just proved, including the modified low product mask. Form the new
$S,T_+,n$ by \eqref{eq:narrow-binary-packing} and define its actual
new $r$. The packed lemma gives
$n^2\mid\binom{2r}{r}$ and $r\ge n^2-1$.

The canonical $g$ and $\ell$ have zero radix-unit digit. Thus $S$
and $T_+$ are even, and the new $r$ is even. Apply the positive
Pell construction of Section~\ref{sec:binary-product} afresh at this
$n,r$: set $U=2^{2r+1}$ and
$Y_P=\lfloor(U+1)^{2r}/U^r\rfloor$, use the main index $2r+1$,
the first index $r+1$, and the second index $L$. The same exact
congruences, ratio gaps and even-index half-parameter construction
make every quotient and every remaining supplied witness positive.
No old packed index or Pell witness is reused.

Both packings still cost the same number of operations. The new
power chain computes $q^2,q^4,n^2=q^8$ in three multiplications,
where the old one computed $q^2,q^4,q^8,n^2=q^{16}$ in four.
The baseline term $q^2$ formerly supplied by
$q^2(1+\theta\lambda)$ is absent, so the reordered mask incurs no
extra addition. This gives 47 multiplications and 42 additions,
proving Theorem~\ref{thm:best89}.

The complete source checker compares all 22 independently written
polynomial residuals with all 89 primitive instructions, including
the same acyclic norm and geometric corrections. The full table is
Appendix~\ref{app:89}. The companion files
\file{BINARY\_PRODUCT\_89\_PROOF.md} and
\file{BASE\_TWO\_PELL\_89\_PROOF.md} state the outer and Pell proofs
separately. Finite regression checks corroborate these arguments;
they do not replace either unbounded positive-domain direction.
